\section{数据来自怎样的生成机制}
\label{sec:11-Mathematical-Statistics/01-Statistical-Models/01}

100 次产品质检，92 次合格。算术告诉你 0.92，但你真正想知道的不是这个数字——你想知道整条生产线的长期合格率大概是多少，想知道明天抽检的合格概率，想知道工艺改进后有没有真的变好。

0.92 是这 100 件产品的总结。它能不能代表更多东西，取决于这 100 件是怎么来的。如果它们是从稳定生产线上随机抽取的，0.92 可能是长期合格率的一个合理近似。如果只抽了白班、只记录了通过初筛的产品、或者每件产品来自不同设备，那 0.92 替谁说话都说不清楚。

数据不会自己解释自己。一张数据表的每一行背后都有一个生成它的过程。统计推断的第一步，不是算均值、画直方图、跑回归，而是问：这些数据是从怎样的机制里产生出来的？

\subsection{总体首先是一种目标分布}

最朴素的想法：有一群我们关心的对象——所有产品、所有患者、所有像素——从中抽出一部分观察，用观察结果推断整体。这群对象叫总体。

在有限调查中，总体可以是一份名单上的全部个体。但在重复实验中，总体是所有可能重复的集合；在时间序列中，总体是随机过程的一条轨线；在机器学习中，总体是图像和标签的联合分布。更一般地，我们把总体理解为我们关心的数据生成分布 \(P\)——一个描述观测如何产生的概率规则。

观测前，样本写成随机变量：
\[
X_1,\ldots,X_n.
\]
大写字母不是排版习惯。\(X_1\) 的意思是：如果我明天重新抽 100 件，第一件的结果可能和今天不同。它是还未发生的可能取值，取决于抽到了哪个个体。

观测后，得到具体数值：
\[
x_1,\ldots,x_n.
\]
小写字母是已经发生的事实。92 就是 92，不会再变。

概率陈述——偏差、方差、覆盖率——作用于随机变量，讨论的是``如果重复抽样会怎样''。实际计算——均值、分位数、p 值——代入观测值。大小写的区别不是在咬文嚼字：它区分了两类不同的讨论，一类关于估计方法在重复中的表现，另一类关于这一次数据给出的具体结论。

但``有一个分布 P''还太模糊。\(P\) 可以是任何东西。如果什么分布都有可能，从 100 个数据点根本无法做有意义的推断——对见过的数据可以精确拟合，对未见过的数据毫无预测力。我们需要缩小候选范围。

\subsection{统计模型是一组候选生成机制}

统计模型做的事就是画一个圈：在这些候选分布里，我们认真考虑数据可能是怎么产生的。圈外的可能性暂时不讨论。

参数模型把这个圈写成：
\[
\mathcal P=\{P_\theta:\theta\in\Theta\}.
\]
\(\theta\) 是未知参数，\(\Theta\) 是参数空间——所有允许的参数值。Bernoulli 模型里：
\[
X_i\sim\operatorname{Bernoulli}(\theta),\qquad0\le\theta\le1,
\]
\(\theta\) 就是合格概率本身。正态模型里 \((\mu,\sigma^2)\) 表示位置和尺度；回归里是系数向量 \(\beta\)；混合模型里还有隐变量的成分权重。

一个模型不保证``真实世界一定就是这些分布中的一个''。它只决定：在这个框架内，哪些推断可以识别，似然函数怎么写，标准误怎么算，结论的适用范围画在哪条线上。选模型就是选一套推理规则。换一个模型，同样的数据可能给出不同的答案——不是因为数据变了，而是因为你在用不同的方式提问。

非参数模型不是``没有假设''。它只是不给分布限定一个低维参数族。你仍然在假定独立性、平滑性、某种尾部行为、数据来自某个抽样设计。这些是实质假设，只是不以两三个参数的形式表达。

\subsection{独立同分布是两条不同假设}

最常见的简写：
\[
X_1,\ldots,X_n\overset{\mathrm{iid}}{\sim}P_\theta
\]
同时包含两层意思：
\begin{itemize}
\tightlist
\item \textbf{同分布}：每个观测来自同一个 \(P_\theta\)。第一个产品和最后一个产品服从完全相同的随机机制，不因为采样顺序而改变分布。
\item \textbf{独立}：知道前面几个样本的结果，不改变你对后面样本的条件分布。出现一个合格品，不意味着下一个更可能合格，也不意味着更不可能。
\end{itemize}

这两条都可能失败，而且失败的方式不同。时间序列中相邻值相关——今天的温度和昨天的温度不独立，但分布可能相同。概念漂移中，后期数据和早期数据的分布不同——同一个传感器，老化后读数的期望值变了。分层抽样中，各组有各自的分布——不是同分布，但在层内可能近似满足。

从有限总体不放回抽样时，观测也不严格独立：抽走一个个体后，剩下的个体被抽中的概率会改变。但如果总体远大于样本量，这种依赖很弱，iid 近似可能足够好。要不要近似、近似到什么程度，必须结合抽样比例和推断目的来判断。

不能把``数据行彼此不同''当成独立的证据。独立性描述的是联合生成机制——数据之间的随机依赖关系——而不是表格里数值是否出现重复。一个全为 1 的 Bernoulli 样本仍然可能是独立的，只是恰好都取到了 1。

回到 92/100。如果我们愿意接受 iid Bernoulli 模型，0.92 就有了一个明确的角色：它是 \(\theta\) 的一个估计，来自一个定义清晰的概率框架。在这个框架内，我们可以讨论它离真 \(\theta\) 有多远，可以构造区间，可以比较不同估计方法的好坏。

模型把 0.92 从``一个数字''变成了``一个参数的影子''。

\subsection{选择偏差不会被大样本消除}

大数定律让样本均值逼近实际抽样分布的期望。但它不会把一个有偏的抽样框变成目标总体。

在线问卷只吸引强烈意见者。即使收到一百万份答卷，你也只能非常精确地估计``愿意填问卷的人''的平均意见，而不是目标人群的平均意见。样本量消除了随机误差，对系统性的选择偏差、测量误差和标签错误完全无效。

统计设计因此先于公式。在动用任何统计方法之前，你需要回答：
\begin{itemize}
\tightlist
\item 谁有机会进入样本？
\item 进入概率是否已知，是否与结果相关？
\item 缺失是随机发生，还是与未观测值相关？
\item 处理是否随机分配？
\item 观测单位之间是否相互干扰？
\end{itemize}
如果这些问题没有答案，置信区间再窄也可能只是在错误中心周围非常自信。

\subsection{可识别性：不同参数能否产生不同分布}

参数 \(\theta\) 可识别，意味着：
\[
P_{\theta_1}=P_{\theta_2}\quad\Longrightarrow\quad\theta_1=\theta_2.
\]
如果两个不同的参数值给出完全相同的观测分布，无论收集多少数据，都不可能把它们区分开。模型本身就不具备区分它们的信息——这不是数据量的问题，是模型设计的问题。

一个典型的例子：只观察正态变量的平方 \(Y=X^2\)，其中 \(X\sim N(\mu,\sigma^2)\)。\(Y\) 的分布只通过 \(\mu^2\) 依赖 \(\mu\)，因此 \((\mu,\sigma^2)\) 和 \((-\mu,\sigma^2)\) 产生完全相同的观测分布。数据中没有任何信息能告诉你 \(\mu\) 的正负号。在神经网络和混合模型中，隐藏单元交换也产生同一分布。

可识别性属于模型和观测机制，不是优化收敛性。算法返回了一个参数值，不代表数据唯一识别了它。如果模型不可识别，再好的优化、再多的数据也无济于事——你得到的只是众多等价答案中的一个。

\subsection{模型错设时结论变成什么}

真实分布 \(P_0\) 可能根本不在你列出的候选族 \(\mathcal P\) 中。此时估计方法仍可能收敛，但它收敛到的不是``真实参数''——真实参数在模型内不存在——而是在某种准则下最接近真实分布的``伪真参数''。

正态均值模型对轻微非正态数据仍可能很好地估计期望——均值的估计对分布假设不太敏感。但基于正态尾部计算的极端概率（比如``超出 6 个标准差的概率''）可能差几个数量级。

模型诊断应针对推断目标：估计均值、估计分位数、估计尾部风险、估计因果效应——每种目标对模型假设的敏感程度不同。不要问``模型正确吗''——所有模型都错。问``在我要做的这件事上，模型的错误会产生多大影响''。
